% GENERATED FILE. Edit canonical lessons, then run npm run generate:books.
% canonical-source-hash: fnv1a64:8dff42f1550654bc
% canonical-lessons: KA-C199-nungu, KA-C199-hasu, KA-C199-madacu, KA-C199-istri-madu, KA-C199-badalisu

\chapter{To swallow, To spread out, To fold, To iron, To change}
\label{ch:ka-to-swallow-to-spread-out-to-fold-to-iron-to-change}
\hlchaptermodality{199}

\begin{chapteropening}
\textbf{By the end of this chapter:} \emph{I can say to swallow, to spread out, to fold, to iron and to change.}

Complete the last lesson of chapter 199: I can say to swallow, to spread out, to fold, to iron and to change.
\end{chapteropening}

\section[\texorpdfstring{\kn{ನುಂಗು} (nuṅgu)}{ನುಂಗು (nuṅgu)}]{\kn{ನುಂಗು} (nuṅgu) — to swallow}

\label{lesson:KA-C199-nungu}



\begin{quote}
Before the new one: say the Kannada for to hold, then the Kannada for to let go.
\end{quote}

\subsection*{You'll want to know: \kn{ನುಂಗು}}
\textbf{\kn{ನುಂಗು}} — \emph{nuṅgu} — \textquotedblleft{}to swallow\textquotedblright{}.

\begin{grammarlens}[title={What you've built}]
One more word for this chapter.
\end{grammarlens}

\subsection*{Your turn}
\begin{itemize}
\raggedright
  \item \emph{Say it:} \emph{nuṅgu}
  \item \emph{Say it:} \emph{nuṅgu}, once more
  \item \emph{Say it:} say \emph{nuṅgu}
  \item \emph{Recall:} say the Kannada for to hold, then the Kannada for to let go, then say \emph{nuṅgu} again
\end{itemize}

\begin{tcolorbox}[breakable,colback=teal!4,colframe=teal!35!black,title={Before you move on}]
What does \kn{ನುಂಗು} mean? (\textquotedblleft{}To swallow\textquotedblright{}.) Say it once more.
\end{tcolorbox}
\section[\texorpdfstring{\kn{ಹಾಸು} (hāsu)}{ಹಾಸು (hāsu)}]{\kn{ಹಾಸು} (hāsu) — to spread out (a mat, bedding)}

\label{lesson:KA-C199-hasu}



\begin{quote}
Before the new one: say the Kannada for to let go, then the Kannada for to swallow.
\end{quote}

\subsection*{You'll want to know: \kn{ಹಾಸು}}
\textbf{\kn{ಹಾಸು}} — \emph{hāsu} — \textquotedblleft{}to spread out (a mat, bedding)\textquotedblright{}.

\begin{grammarlens}[title={What you've built}]
One more word for this chapter.
\end{grammarlens}

\subsection*{Your turn}
\begin{itemize}
\raggedright
  \item \emph{Say it:} \emph{hāsu}
  \item \emph{Say it:} \emph{hāsu}, once more
  \item \emph{Say it:} say \emph{hāsu}
  \item \emph{Recall:} say the Kannada for to let go, then the Kannada for to swallow, then say \emph{hāsu} again
\end{itemize}

\begin{tcolorbox}[breakable,colback=teal!4,colframe=teal!35!black,title={Before you move on}]
What does \kn{ಹಾಸು} mean? (\textquotedblleft{}To spread out (a mat, bedding)\textquotedblright{}.) Say it once more.
\end{tcolorbox}
\section[\texorpdfstring{\kn{ಮಡಚು} (maḍacu)}{ಮಡಚು (maḍacu)}]{\kn{ಮಡಚು} (maḍacu) — to fold}

\label{lesson:KA-C199-madacu}



\begin{quote}
Before the new one: say the Kannada for to swallow, then the Kannada for to spread out.
\end{quote}

\subsection*{You'll want to know: \kn{ಮಡಚು}}
\textbf{\kn{ಮಡಚು}} — \emph{maḍacu} — \textquotedblleft{}to fold\textquotedblright{}.

\begin{grammarlens}[title={What you've built}]
One more word for this chapter.
\end{grammarlens}

\subsection*{Your turn}
\begin{itemize}
\raggedright
  \item \emph{Say it:} \emph{maḍacu}
  \item \emph{Say it:} \emph{maḍacu}, once more
  \item \emph{Say it:} say \emph{maḍacu}
  \item \emph{Recall:} say the Kannada for to swallow, then the Kannada for to spread out, then say \emph{maḍacu} again
\end{itemize}

\begin{tcolorbox}[breakable,colback=teal!4,colframe=teal!35!black,title={Before you move on}]
What does \kn{ಮಡಚು} mean? (\textquotedblleft{}To fold\textquotedblright{}.) Say it once more.
\end{tcolorbox}
\section[\texorpdfstring{\kn{ಇಸ್ತ್ರಿ} \kn{ಮಾಡು} (istri māḍu)}{ಇಸ್ತ್ರಿ ಮಾಡು (istri māḍu)}]{\kn{ಇಸ್ತ್ರಿ} \kn{ಮಾಡು} (istri māḍu) — to iron (clothes)}

\label{lesson:KA-C199-istri-madu}



\begin{quote}
Before the new one: say the Kannada for to spread out, then the Kannada for to fold.
\end{quote}

\subsection*{You'll want to know: \kn{ಇಸ್ತ್ರಿ} \kn{ಮಾಡು}}
\textbf{\kn{ಇಸ್ತ್ರಿ} \kn{ಮಾಡು}} — \emph{istri māḍu} — \textquotedblleft{}to iron (clothes)\textquotedblright{}.

\begin{grammarlens}[title={What you've built}]
One more word for this chapter.
\end{grammarlens}

\subsection*{Your turn}
\begin{itemize}
\raggedright
  \item \emph{Say it:} \emph{istri māḍu}
  \item \emph{Say it:} \emph{istri māḍu}, once more
  \item \emph{Say it:} say \emph{istri māḍu}
  \item \emph{Recall:} say the Kannada for to spread out, then the Kannada for to fold, then say \emph{istri māḍu} again
\end{itemize}

\begin{tcolorbox}[breakable,colback=teal!4,colframe=teal!35!black,title={Before you move on}]
What does \kn{ಇಸ್ತ್ರಿ} \kn{ಮಾಡು} mean? (\textquotedblleft{}To iron (clothes)\textquotedblright{}.) Say it once more.
\end{tcolorbox}
\section[\texorpdfstring{\kn{ಬದಲಿಸು} (badalisu)}{ಬದಲಿಸು (badalisu)}]{\kn{ಬದಲಿಸು} (badalisu) — to change}

\label{lesson:KA-C199-badalisu}



\begin{quote}
Before the new one: say the Kannada for to fold, then the Kannada for to iron.
\end{quote}

\subsection*{You'll want to know: \kn{ಬದಲಿಸು}}
\textbf{\kn{ಬದಲಿಸು}} — \emph{badalisu} — \textquotedblleft{}to change\textquotedblright{}.

\begin{grammarlens}[title={What you've built}]
One more word for this chapter.
\end{grammarlens}

\subsection*{Your turn}
\begin{itemize}
\raggedright
  \item \emph{Say it:} \emph{badalisu}
  \item \emph{Say it:} \emph{badalisu}, once more
  \item \emph{Say it:} say \emph{badalisu}
  \item \emph{Recall:} say the Kannada for to fold, then the Kannada for to iron, then say \emph{badalisu} again
\end{itemize}

\begin{tcolorbox}[breakable,colback=teal!4,colframe=teal!35!black,title={Before you move on}]
What does \kn{ಬದಲಿಸು} mean? (\textquotedblleft{}To change\textquotedblright{}.) Say it once more.
\end{tcolorbox}
